When a gas leak occurs in a plant, a storage facility or a confined space, fixed sensor networks answer the question \textit{``is there a leak''} well, but they answer the question \textit{``where is the leak''} poorly. The second question is still largely solved by sending a human with a handheld detector into the hazardous area.

This project, \textbf{GasSeeker}, builds an autonomous ground robot carrying \textbf{a single gas sensor} to localise a gas source inside a 2 $\times$ 2 m test arena, together with a LoRa ground station so that the operator can stay outside the hazardous area. The single-sensor constraint is not a cost-cutting measure but \textbf{the condition that defines the problem}: when instantaneous spatial concentration differences cannot be measured, the robot must substitute time, motion and wind-direction information. The report gives the physical argument for this constraint, grounded in the published literature on the intermittent structure of turbulent odour plumes and in a signal-to-noise analysis showing that increasing the separation between two on-board sensors does not improve directional discrimination.

Three source-localisation algorithms --- exhaustive grid search, concentration-gradient following, and moth-inspired surge-casting --- are implemented on \textbf{the same hardware platform, the same measurement procedure and the same scoring criterion}, and compared under two flow regimes: pure diffusion and intermittent wind-driven dispersion. All algorithm code lives in a hardware-independent C++ layer, so \textbf{the exact code that will be flashed onto the microcontroller is the code that runs in simulation}.

Over 30 simulated runs per algorithm $\times$ environment cell (180 runs in total), the results show a clear and statistically significant reversal: gradient following achieves a 67 \% success rate under diffusion but drops to 17 \% under wind (\textit{p} = 0.0002); surge-casting moves the opposite way, from 17 \% to 57 \% (\textit{p} = 0.0028); exhaustive search remains statistically unchanged at 77 \% and 63 \% (\textit{p} = 0.40) while taking roughly 2.2 to 2.7 times longer. The conclusion is \textbf{not} that one algorithm is best, but a \textbf{selection rule keyed to the flow regime}.

The most important limitation is stated explicitly: \textbf{all comparative figures are simulation results, not measurements on the physical robot}. The hardware has been designed, wired and programmed, and the firmware compiles cleanly, but the experimental campaign could not be completed within the available time.

\needspace{6\baselineskip}\noindent \textbf{Keywords:} gas source localisation, robotic olfaction, MQ-3 sensor, surge-casting, gradient following, turbulent plume, ESP32-S3, LoRa.
